\section{Generación de reportes}
\label{sec:met_generacion}

% Extensión estimada: 3--3.5 páginas, con dos tablas (configuración de LoRA y del entrenamiento; decodificación).
% Fuentes: Codigo/outputs/checkpoints/stage2_mini_vlm_v2 (lora_config.json y best_model.pt, módulos y épocas);
% Codigo/slurm/mendieta/v2/train_vlm_s2*.sh (tasa, lote, LoRA); Codigo/src/training/two_stage.py (receta de ambas
% etapas, freeze, instrucción de entrenamiento); Codigo/scripts/train_vlm_stage1.py y train_vlm_stage2.py;
% Codigo/scripts/generate_report.py (instrucción de inferencia); Codigo/slurm/mendieta/v3/ y v4/ (384 tokens, semilla);
% Codigo/slurm/mendieta/pilot_ensemble_v2_llava.sh (motivo de 384); Codigo/data/splits/skincare_*_splits.json;
% conteo propio sobre Codigo/data/raw/skincare/SkinCoT/EN; Codigo/docs/resul/analisis_pipeline_avanzado_resul_20260908.md
% (penalización de repetición). Citas: \cite{shen2025skincare}, \cite{zeng2025mmskin}, \cite{abdin2024phi3},
% \cite{liu2023visual}, \cite{li2023llavamed}, \cite{jiang2023mistral}, \cite{hu2021lora}, \cite{tschandl2018ham10000}.

\subsection{Partes del reporte que genera el trabajo}

El generador tiene que producir las dos partes del reporte que la Sección~\ref{sec:met_tareas} define para la tarea
de generación: la descripción de lo que se ve en la imagen y la impresión diagnóstica. El diagnóstico diferencial y
la recomendación de manejo quedan fuera. En la Sección~\ref{sec:met_clasif_vlm} los \glspl{vlm} clasifican, es decir,
eligen una etiqueta de un inventario cerrado. Redactar un texto libre con la estructura de un reporte exige otra forma
de salida, y por eso los modelos se ajustan para generar.
\pendiente{dato: en tres o cuatro reportes de una corrida, verificar si la impresión diagnóstica aparece en la
condición base, sin diagnóstico declarado (Q15). Si no aparece, decirlo aquí y remitir a la Sección~\ref{sec:met_pipeline}.}

\subsection{Generadores}

Los dos \glspl{vlm} generales de la Sección~\ref{sec:met_clasif_vlm} se ajustan para generar. Mini-VLM es
\texttt{microsoft/Phi-3.5-vision-instruct}~\cite{abdin2024phi3}, y LLaVA-Med v1.5 es un modelo de la familia
LLaVA~\cite{liu2023visual} adaptado a la biomedicina~\cite{li2023llavamed}, con el modelo de lenguaje Mistral-7B
\cite{jiang2023mistral}. Son los mismos modelos base que se usan para clasificar, así que el ajuste es lo que cambia su
salida de una etiqueta a un texto. Ambos son de propósito general, y falta un generador especializado en dermatología
con el cual compararlos.

\textcite{zeng2025mmskin} presentan SkinVL junto con MM-Skin, un conjunto de pares imagen--texto de dermatología.
SkinVL es LLaVA-Med v1.5 con Mistral-7B ajustado a dermatología con MM-Skin.
\pendiente{dato: el Capítulo~2 atribuye a zeng2025mmskin un ajuste con proyección y adaptadores LoRA, mientras que la
inspección del repositorio dice que los pesos publicados (\texttt{zwq803/SkinVL-MM}) son completos y sin adaptadores
(Q11). Describir el método según el artículo y anotar la inspección como dato propio.}
Los pesos publicados están en el formato original de LLaVA-Med y no incluyen el codificador visual, así que se
convirtieron al formato de Hugging Face: los pesos del modelo de lenguaje y del proyector se trasplantaron sobre la base
\texttt{chaoyinshe/llava-med-v1.5-mistral-7b-hf}, que aporta el codificador visual. No se ajustó de nuevo. Se incluye
como generador especializado de comparación, para saber si la especialización en dermatología cambia el texto frente
a los generales.

\subsection{Ajuste en dos etapas sobre SkinCaRe}

El ajuste tiene dos etapas, y en ambas la instrucción es la misma: «Describe this skin lesion.». En la primera se entrena
solo el proyector visual, la capa que traduce las características de la imagen al espacio del modelo de lenguaje, con
pares imagen--caption de \emph{SkinCAP}. El codificador visual y el modelo de lenguaje quedan congelados, de modo que
esta etapa solo alinea la imagen con el texto. En la segunda se aplica \gls{lora} sobre los pares de \emph{SkinCoT} en
inglés. Los adaptadores se colocan en los módulos llamados \texttt{q\_proj} y \texttt{v\_proj}, y ese filtro por nombre
alcanza partes distintas de cada modelo. En LLaVA-Med coincide con las proyecciones de atención del modelo de lenguaje y
también con las del codificador visual \gls{clip}. En Mini-VLM, cuyo modelo de lenguaje fusiona las tres proyecciones de
atención en un único módulo (\texttt{qkv\_proj}), el filtro solo coincide con el codificador visual: el adaptador quedó
en las 24 capas de ese codificador y el modelo de lenguaje no se ajustó.
\pendiente{decisión del autor: declarar así el ajuste de Mini-VLM o reentrenarlo con LoRA sobre \texttt{qkv\_proj}.
Verificar los módulos del checkpoint v2 de LLaVA-Med, que no está en la copia local (en la v1: 64 módulos en el modelo
de lenguaje y 48 en el codificador visual), y si el proyector se entrena en la etapa 2 (Q4).}

El ajuste no usa \emph{HAM10000}, que trae una etiqueta diagnóstica por imagen y no un texto, así que no sirve como
referencia para entrenar el generador~\cite{tschandl2018ham10000}. Su partición de prueba, de 1\,503 imágenes en la partición por lesión, se reserva
para medir la fidelidad diagnóstica (Sección~\ref{sec:met_protocolo}). \emph{SkinCaRe} trae, en cambio, un texto de
clínicos para cada imagen: las descripciones de \emph{SkinCAP}, escritas por dermatólogos, y el razonamiento diagnóstico
revisado de \emph{SkinCoT}~\cite{shen2025skincare}. De los 4\,000 pares de \emph{SkinCAP}, la partición usa 3\,561.
\pendiente{dato: motivo de la exclusión de 439 pares (Q10)} \emph{SkinCoT} aporta 3\,041 textos, y la segunda etapa
ve 23 clases de enfermedades, no las siete de \emph{HAM10000}.
\pendiente{dato: confirmar qué anota cada subconjunto y si la impresión diagnóstica está en el texto de SkinCoT o se
obtiene de su conclusión (Q14).}
Sin embargo, los textos de \emph{SkinCoT} son razonamientos y no reportes, lo que tiene consecuencias para el
entrenamiento.

Por la restricción de cómputo de la Sección~\ref{sec:met_tareas}, los modelos se adaptan
con \gls{lora} en lugar de ajustar todos sus parámetros. La técnica congela los pesos originales y entrena matrices de
bajo rango que se suman a ellos~\cite{hu2021lora}, lo que reduce el número de parámetros entrenados
(Sección~\ref{sec:marco_llm_lora}). La Tabla~\ref{tab:met_lora} fija, para cada etapa, los valores del entrenamiento:
rango, escala, dropout y módulos de LoRA, tasa de aprendizaje, optimizador, épocas y tamaño de lote.

\begin{table}[H]
    \centering
    \caption{Configuración de LoRA y del entrenamiento en las dos etapas del ajuste.}
    \label{tab:met_lora}
    \small
    \begin{tabularx}{\textwidth}{@{}>{\raggedright\arraybackslash}p{4.2cm}
        >{\raggedright\arraybackslash}X@{}}
        \toprule
        \textbf{Parámetro} & \textbf{Valor} \\
        \midrule
        \multicolumn{2}{@{}l}{\emph{Etapa 1: proyector visual}} \\
        \midrule
        Parámetros entrenables & Proyector visual (LLaVA-Med: \texttt{multi\_modal\_projector}; Mini-VLM: \texttt{img\_projection}, \texttt{glb\_GN} y \texttt{sub\_GN}) \\
        \filasep
        Congelados & Codificador visual y modelo de lenguaje \\
        \filasep
        Datos & 2\,492 / 534 / 535 (entrenamiento / validación / prueba) \\
        \filasep
        Optimizador & AdamW, decaimiento 0, eps $10^{-7}$ \\
        \filasep
        Planificador & Coseno \\
        \filasep
        Tasa, épocas, lote y paciencia & a verificar \\
        \midrule
        \multicolumn{2}{@{}l}{\emph{Etapa 2: LoRA}} \\
        \midrule
        Rango ($r$) & 64 \\
        \filasep
        Escala ($\alpha$) & 64 \\
        \filasep
        Dropout de LoRA & 0.0 \\
        \filasep
        Módulos & \texttt{q\_proj}, \texttt{v\_proj} (ver el texto) \\
        \filasep
        Proyector & a verificar \\
        \filasep
        Datos & 2\,128 / 456 / 457 (entrenamiento / validación / prueba) \\
        \filasep
        Instrucción & «Describe this skin lesion.» \\
        \filasep
        Tasa de aprendizaje & $2.495 \times 10^{-5}$ \\
        \filasep
        Optimizador & Decaimiento 0.01, eps $10^{-7}$ \\
        \filasep
        Épocas & 10, paciencia 3 \\
        \filasep
        Lote y acumulación & 4 y 8 (32 efectivos) \\
        \filasep
        Longitud máxima & 2\,048 \\
        \filasep
        Recorte de gradiente & 1.0 \\
        \filasep
        Precisión & bf16 en GPU de arquitectura Ampere o posterior \\
        \filasep
        Cuantización & a verificar \\
        \filasep
        Semilla & 42 \\
        \bottomrule
    \end{tabularx}
\end{table}
\pendiente{Tabla~\ref{tab:met_lora}: completar la tasa, las épocas, el lote y la paciencia de la etapa 1 (Q1), si el
proyector se entrena en la etapa 2 (Q4), el lote real según el log (Q6) y si hubo cuantización (Q7).}

\subsection{Reporte y razonamiento}

En \emph{SkinCoT} el texto es una narrativa de razonamiento diagnóstico, y sus pasos van entre bloques
\texttt{\textless think\textgreater} y \texttt{\textless /think\textgreater}. En 2\,942 de los 3\,041 textos aparece al
menos uno (96.74\,\%). El entrenamiento usa el texto tal cual, así que el modelo aprende también a reproducir esos
bloques, y su salida mezcla razonamiento y reporte. Los bloques se conservan en el entrenamiento porque el razonamiento
es auditable: permite ver cómo el modelo llega a su descripción. Ese mismo texto no sirve para medir la calidad del
reporte, y por eso la separación se hace en la evaluación.

La instrucción de inferencia es la del entrenamiento, «Describe this skin lesion.». Un prompt de inferencia distinto
cambiaría lo que recibe el modelo respecto de lo que aprendió, así que se conserva la misma. Para medir, los bloques de
razonamiento se quitan del texto generado y las métricas de generación se calculan solo sobre el reporte, con el
criterio de la Sección~\ref{sec:met_protocolo}. Si se midiera el texto completo, las menciones de diagnósticos que el
razonamiento solo hipotetiza contarían como aciertos. La magnitud de ese efecto se presenta en el
Capítulo~\ref{cp:res_gen}.

\subsection{Decodificación}

La generación es voraz (\emph{greedy decoding}, \texttt{do\_sample=False}): en cada paso elige el token más probable, en
lugar de muestrear como en la Sección~\ref{sec:marco_llm_grandes}. Cada imagen se procesa en una llamada, con un máximo
de 384 tokens nuevos y semilla 42 (Tabla~\ref{tab:met_decodificacion}). La decodificación voraz tiende a repetirse en la
generación abierta, y el trabajo la acepta por reproducibilidad. El límite de 384 tokens reemplazó al de 256 del
piloto, para reducir el truncamiento antes de que aparezca el diagnóstico. No se usó
penalización de repetición ni \texttt{no\_repeat\_ngram\_size} en las corridas finales: en Mini-VLM la penalización hace
fallar la generación, y el posprocesamiento que quita repeticiones se evaluó y no se adopta, con su resultado en el
Capítulo~\ref{cp:res_gen}. Con esto el generador redacta el texto, pero no decide si la lesión es lo bastante clara para
afirmar un diagnóstico. Esa decisión, y la evidencia que la sostiene, es la que incorpora el \emph{pipeline} de la
Sección~\ref{sec:met_pipeline}.

\begin{table}[H]
    \centering
    \caption{Decodificación de los tres generadores en las corridas finales.}
    \label{tab:met_decodificacion}
    \small
    \begin{tabularx}{\textwidth}{@{}>{\raggedright\arraybackslash}p{4.4cm} >{\raggedright\arraybackslash}X@{}}
        \toprule
        \textbf{Parámetro} & \textbf{Valor} \\
        \midrule
        Estrategia & Voraz (\texttt{do\_sample=False}) \\
        \filasep
        Máximo de tokens nuevos & 384 en las corridas finales; 512 en las dos primeras iteraciones \\
        \filasep
        Llamadas & Una imagen por llamada \\
        \filasep
        Penalización de repetición & No se usa en las corridas finales \\
        \filasep
        Instrucción de inferencia & «Describe this skin lesion.», la misma del entrenamiento \\
        \filasep
        Semilla & 42 \\
        \filasep
        Posprocesamiento & Se quitan los bloques \texttt{\textless think\textgreater} antes de medir \\
        \bottomrule
    \end{tabularx}
\end{table}
